\section{Memory, artifacts, and distributed knowledge}
\label{sec:memory}
\subsection{Storage is not collective recall}
A swarm's persistent state can exceed any participant's context window. That observation does not imply that the system can retrieve, reconcile, or use all of it. Three quantities must be separated: information stored, information accessible within a resource budget, and task-relevant information correctly integrated into an answer. Summing context-window lengths measures none of the latter two. It double-counts repeated content, ignores inaccessible private facts, and says little about the cost of reading shared state.

Single-agent systems are essential controls. Recursive Language Models place a large input in a programmable environment and let a root process inspect it and call models on selected fragments. The reported accuracy is 91.3\% on a BrowseComp-Plus setting containing 6--11 million tokens, beyond the root model's native window \citep{recursive-language-models}. Most subcalls in the main experiments are sequential and do not communicate. File-based coding agents likewise answer long-context questions without a simultaneously interacting team \citep{coding-agents-long-context}. These results demonstrate externalized processing, not an intrinsic multi-agent memory advantage.

Context management can also account for an apparent swarm benefit. The Kimi comparison changes substantially when the single agent can clean its context, while a context-as-file approach shows that changing how state is edited can matter within both single-agent and multi-agent systems \citep{kimi-agent-swarm,context-language-models}. A claim about collective memory should therefore compare against a strong individual with the same retrieval tools, external storage, and opportunity to revisit evidence.

\subsection{Shared artifacts carry both discoveries and errors}
The corpus contains several different memory mechanisms. A Git history preserves an artifact and its ancestry; a task board records ownership; a claim graph records dependencies; a note archive supports retrieval; a simulated environment preserves objects that later agents can inspect. These stores differ in what can be overwritten and in how a reader knows whether an entry is current or verified. Calling all of them ``shared memory'' is useful only if the update and retrieval rules remain visible.

The Fermat formalization campaign uses a shared graph of claims to maintain a project across many agent sessions \citep{fermat-last-theorem}. The engineering account associates this structure with recovery from earlier attempts that lost project state. Formal checking constrains the acceptance of proofs, but the campaign does not isolate the graph's effect from model capability and the rest of the platform. Agora preserves contributions in a Git graph, allowing successors to inherit both results and failed attempts; much of its reported optimization gain occurs early, followed by concentration on one branch \citep{agora-git-memory}. Provenance supports reuse while also making a successful lineage an attractor.

CORAL offers a more direct ablation of accumulated notes: removing them worsens the single agent's kernel optimization result, while a shared multi-agent run builds many record improvements on others' commits \citep{coral}. The same system still requires an evaluator that accepts or rejects executable artifacts. Its evidence is strongest for this class of search, where a concrete candidate and a reproducible score can travel together. It does not imply equally reliable transfer of unverified scientific interpretations.

The enzyme-discovery campaign combines version history with a curator that adds results to a shared knowledge base \citep{anthropic-art-enzymes}. The reported discovery depends on one worker reading raw DNA and noticing a feature not explicitly requested; ten restarts do not reproduce that observation. This is a useful counterweight to a simple cumulative-memory story. A system can store abundant summaries and still miss a detail that is absent from them. The outcome is evidence about a particular campaign and its information path, not a measured probability that adding workers will recover such details.

\subsection{Knowledge can be present and still remain unused}
HiddenBench distributes shared and private clues among agents. The team averages 30.1\% accuracy, versus 80.7\% for one agent given all the facts. Requiring disclosure of facts early in the protocol greatly improves some model configurations \citep{hiddenbench}. The important failure is not necessarily reasoning over a known set of clues: it is failing to elicit the clues before deciding. A single-agent full-information baseline diagnoses that gap but also gives the single agent more information at the outset, which must remain explicit.

Silo-Bench increases the number of agents holding distributed data up to 100. Performance falls with group size in the reported DeepSeek configuration, and premature submission or nonconvergence accounts for many examined failures \citep{silo-bench}. Agents can choose a plausible graph without successfully completing the data exchange. This separates organizational appearance from information-flow completion. A diagram that resembles a distributed algorithm is not evidence that its invariants hold.

The Opus 5.5 knowledge-building task measures a downstream reader answering hidden questions from the notes produced in 24 hours. Scores rise from 0.53 for one agent to 0.70 for ten and 0.74 for one hundred \citep{opus55-card}. These are accuracy outcomes at one deadline. The smaller marginal improvement above ten agents neither measures a memory-capacity ceiling nor proves saturation across tasks. The unreported token budget and the question distribution both matter. Similarly, the slight advantage of private over shared memory in one LLMA-Mem ablation is a local finding, not a general case against common stores \citep{scaling-teams-or-time}.

\subsection{Population memory without live collaboration}
Transactive-memory retrieval can reuse trajectories from a population of models even when the current agent acts alone \citep{multi-agent-transactive-memory}. Its gains test the value of an experience archive and a retrieval policy, not the effect of simultaneous communication. The composition of that archive is itself a confound: one abundant source can dominate both quantity and utility. Scaling the index does not necessarily scale the diversity of evidence.

SwarmWorld provides an environmental counterpart. Agents leave executable constructions in a shared world, and later agents mostly encounter reusable technologies through objects rather than messages. The collective improves a portfolio-level measure while the strongest individual artifact can still come from isolated search \citep{prev-swarmworld}. This disagreement between portfolio and maximum quality is informative. A memory system can improve robustness or breadth without producing the single best candidate, so the outcome must match the intended use.

\subsection{A memory evaluation should test the whole path}
The relevant path is write, retain, retrieve, verify, reconcile, and act. Evaluation should inject known facts with timestamps and provenance, then measure whether they survive updates and become available to the agents that need them. Contradictory claims, obsolete instructions, and concurrent edits test different failure modes. The governed-memory work identifies leakage, staleness, persistent contradictions, and lost provenance as distinct problems \citep{governed-shared-memory}.

Such tests should compare private, shared, and hybrid stores under equal read and write budgets. Report memory growth and access latency alongside downstream accuracy. A small correct index can outperform a much larger archive that overwhelms readers; a store that preserves every message can also preserve malicious instructions. Collective memory is therefore an allocation and trust problem as well as a storage problem, linking this section to both organizational design and safety.
